% GENERATED FILE. Edit canonical lessons, then run npm run generate:books.
% canonical-source-hash: fnv1a64:f22f8e224dfb65a6
% canonical-lessons: BN-C243-bomi, BN-C243-chulkani, BN-C243-dari, BN-C243-goph, BN-C243-chibuk

\chapter{Vomiting, Itch, Beard, Moustache, Chin}
\label{ch:bn-vomiting-itch-beard-moustache-chin}
\hlchaptermodality{243}

\begin{chapteropening}
\textbf{By the end of this chapter:} \emph{I can say vomiting, an itch, a beard, a moustache and a chin.}

Complete the last lesson of chapter 243: I can say vomiting, an itch, a beard, a moustache and a chin.
\end{chapteropening}

\section[\texorpdfstring{\bn{বমি} (bômi)}{বমি (bômi)}]{\bn{বমি} (bômi) — vomiting}

\label{lesson:BN-C243-bomi}



\begin{quote}
Before the new one: say the Bengali for to wrap, then the Bengali for to dig.
\end{quote}

\subsection*{You'll want to know: \bn{বমি}}
\textbf{\bn{বমি}} — \emph{bômi} — \textquotedblleft{}vomiting\textquotedblright{}.

\begin{cousinweb}
\textbf{\bn{বমি}} (\emph{bômi}) is Sanskrit \emph{vami}, \textquotedblleft{}vomiting\textquotedblright{}, taken over whole as a learned loan, a tatsama.
\end{cousinweb}

\begin{grammarlens}[title={What you've built}]
One more word for this chapter.
\end{grammarlens}

\subsection*{Your turn}
\begin{itemize}
\raggedright
  \item \emph{Say it:} \emph{bômi}
  \item \emph{Say it:} \emph{bômi}, once more
  \item \emph{Say it:} say \emph{bômi}
  \item \emph{Recall:} say the Bengali for to wrap, then the Bengali for to dig, then say \emph{bômi} again
\end{itemize}

\begin{tcolorbox}[breakable,colback=teal!4,colframe=teal!35!black,title={Before you move on}]
What does \bn{বমি} mean? (\textquotedblleft{}Vomiting\textquotedblright{}.) Say it once more.
\end{tcolorbox}
\section[\texorpdfstring{\bn{চুলকানি} (chulkāni)}{চুলকানি (chulkāni)}]{\bn{চুলকানি} (chulkāni) — an itch}

\label{lesson:BN-C243-chulkani}



\begin{quote}
Before the new one: say the Bengali for to dig, then the Bengali for vomiting.
\end{quote}

\subsection*{You'll want to know: \bn{চুলকানি}}
\textbf{\bn{চুলকানি}} — \emph{chulkāni} — \textquotedblleft{}an itch\textquotedblright{}.

\begin{cousinweb}
\textbf{\bn{চুলকানি}} (\emph{chulkāni}) is the noun of \emph{chulkāno}, \textquotedblleft{}to scratch\textquotedblright{}, whose origin is not settled: CDIAL has no entry for it.
\end{cousinweb}

\begin{grammarlens}[title={What you've built}]
One more word for this chapter.
\end{grammarlens}

\subsection*{Your turn}
\begin{itemize}
\raggedright
  \item \emph{Say it:} \emph{chulkāni}
  \item \emph{Say it:} \emph{chulkāni}, once more
  \item \emph{Say it:} say \emph{chulkāni}
  \item \emph{Recall:} say the Bengali for to dig, then the Bengali for vomiting, then say \emph{chulkāni} again
\end{itemize}

\begin{tcolorbox}[breakable,colback=teal!4,colframe=teal!35!black,title={Before you move on}]
What does \bn{চুলকানি} mean? (\textquotedblleft{}An itch\textquotedblright{}.) Say it once more.
\end{tcolorbox}
\section[\texorpdfstring{\bn{দাড়ি} (dāṛi)}{দাড়ি (dāṛi)}]{\bn{দাড়ি} (dāṛi) — a beard}

\label{lesson:BN-C243-dari}



\begin{quote}
Before the new one: say the Bengali for vomiting, then the Bengali for an itch.
\end{quote}

\subsection*{You'll want to know: \bn{দাড়ি}}
\textbf{\bn{দাড়ি}} — \emph{dāṛi} — \textquotedblleft{}a beard\textquotedblright{}.

\begin{cousinweb}
\textbf{\bn{দাড়ি}} (\emph{dāṛi}) is a tadbhava, a word inherited and reshaped by sound change. It comes from \emph{*dāṁṣṭra}, \textquotedblleft{}fang, tusk, beard\textquotedblright{}, a form CDIAL reconstructs from its descendants (CDIAL 6250). CDIAL lists Bengali \emph{dāṛi}, \textquotedblleft{}beard, chin\textquotedblright{}.
\end{cousinweb}

\begin{grammarlens}[title={What you've built}]
One more word for this chapter.
\end{grammarlens}

\subsection*{Your turn}
\begin{itemize}
\raggedright
  \item \emph{Say it:} \emph{dāṛi}
  \item \emph{Say it:} \emph{dāṛi}, once more
  \item \emph{Say it:} say \emph{dāṛi}
  \item \emph{Recall:} say the Bengali for vomiting, then the Bengali for an itch, then say \emph{dāṛi} again
\end{itemize}

\begin{tcolorbox}[breakable,colback=teal!4,colframe=teal!35!black,title={Before you move on}]
What does \bn{দাড়ি} mean? (\textquotedblleft{}A beard\textquotedblright{}.) Say it once more.
\end{tcolorbox}
\section[\texorpdfstring{\bn{গোঁফ} (gõph)}{গোঁফ (gõph)}]{\bn{গোঁফ} (gõph) — a moustache}

\label{lesson:BN-C243-goph}



\begin{quote}
Before the new one: say the Bengali for an itch, then the Bengali for a beard.
\end{quote}

\subsection*{You'll want to know: \bn{গোঁফ}}
\textbf{\bn{গোঁফ}} — \emph{gõph} — \textquotedblleft{}a moustache\textquotedblright{}.

\begin{cousinweb}
\textbf{\bn{গোঁফ}} (\emph{gõph}) is a tadbhava, a word inherited and reshaped by sound change. It comes from \emph{*guppha}, \textquotedblleft{}something strung together\textquotedblright{}, a form CDIAL reconstructs from its descendants (CDIAL 4203). CDIAL lists Bengali \emph{gõph}, \textquotedblleft{}moustache\textquotedblright{}, beside Sanskrit \emph{gumpha}, \textquotedblleft{}a whisker\textquotedblright{}.
\end{cousinweb}

\begin{grammarlens}[title={What you've built}]
One more word for this chapter.
\end{grammarlens}

\subsection*{Your turn}
\begin{itemize}
\raggedright
  \item \emph{Say it:} \emph{gõph}
  \item \emph{Say it:} \emph{gõph}, once more
  \item \emph{Say it:} say \emph{gõph}
  \item \emph{Recall:} say the Bengali for an itch, then the Bengali for a beard, then say \emph{gõph} again
\end{itemize}

\begin{tcolorbox}[breakable,colback=teal!4,colframe=teal!35!black,title={Before you move on}]
What does \bn{গোঁফ} mean? (\textquotedblleft{}A moustache\textquotedblright{}.) Say it once more.
\end{tcolorbox}
\section[\texorpdfstring{\bn{চিবুক} (chibuk)}{চিবুক (chibuk)}]{\bn{চিবুক} (chibuk) — a chin}

\label{lesson:BN-C243-chibuk}



\begin{quote}
Before the new one: say the Bengali for a beard, then the Bengali for a moustache.
\end{quote}

\subsection*{You'll want to know: \bn{চিবুক}}
\textbf{\bn{চিবুক}} — \emph{chibuk} — \textquotedblleft{}a chin\textquotedblright{}.

\begin{cousinweb}
\textbf{\bn{চিবুক}} (\emph{chibuk}) is Sanskrit \emph{cibuka}, \textquotedblleft{}chin\textquotedblright{}, taken over whole as a learned loan, a tatsama.
\end{cousinweb}

\begin{grammarlens}[title={What you've built}]
One more word for this chapter.
\end{grammarlens}

\subsection*{Your turn}
\begin{itemize}
\raggedright
  \item \emph{Say it:} \emph{chibuk}
  \item \emph{Say it:} \emph{chibuk}, once more
  \item \emph{Say it:} say \emph{chibuk}
  \item \emph{Recall:} say the Bengali for a beard, then the Bengali for a moustache, then say \emph{chibuk} again
\end{itemize}

\begin{tcolorbox}[breakable,colback=teal!4,colframe=teal!35!black,title={Before you move on}]
What does \bn{চিবুক} mean? (\textquotedblleft{}A chin\textquotedblright{}.) Say it once more.
\end{tcolorbox}
